\documentclass[11pt,a4paper]{article}
\usepackage[utf8]{inputenc}
\usepackage[margin=2.5cm]{geometry}
\usepackage{amsmath,amssymb,bm}
\usepackage{booktabs}
\usepackage{tabularx}
\usepackage{enumitem}
\usepackage[hidelinks]{hyperref}

\newcommand{\vect}[1]{\bm{#1}}
\newcommand{\rr}{\vect{r}}
\newcommand{\om}{\vect{\omega}}
\newcommand{\uu}{\vect{u}}
\newcommand{\nn}{\hat{\vect{n}}}
\newcommand{\Myr}{\,\mathrm{Myr}}
\newcommand{\km}{\,\mathrm{km}}
\newcommand{\gs}{\nabla_{\!s}}

\title{A Kinematic Plate Tectonics Model\\for Procedural World Generation}
\author{Design notes}
\date{November 2026}

\begin{document}
\maketitle

\begin{abstract}
\noindent This document specifies the largest-scale stage of a multi-scale
world generator: a plate tectonics simulation on a full sphere. Plates are
rigid caps that rotate exactly, crust lives in each plate's own frame so it is
never smeared by advection, a low-resolution mantle temperature field supplies
basal drag, and rifting follows weak lines under tension. Elevation is a
diagnostic, not a state variable. The model is intended for a game, so
plausibility and robustness to coarse resolution take priority over accuracy.
Section~\ref{sec:status} separates what is established geophysics from what is
untested construction.
\end{abstract}

\section{Design principles}

\begin{enumerate}[leftmargin=*]
\item \textbf{Plates are given, not emergent.} Obtaining plates from mantle
convection requires delicate rheology and high resolution. Here plates are
explicit objects.
\item \textbf{Exact updates wherever possible.} Rigid rotation and linear
relaxation both have closed-form updates that hold for any time step
$\Delta t$.
\item \textbf{Rates, not events.} Anything irreversible (rifting, welding,
hotspot birth) is expressed as a rate per unit time, so its statistics do not
depend on $\Delta t$.
\item \textbf{Geometric conservation.} Crust is created and destroyed in
proportion to the area swept by plate motion, which is independent of grid
resolution.
\item \textbf{Carry only what feeds back.} The tectonic stage stores the
fields that influence plate dynamics (crust type, thickness, age, damage),
plus history fields that the next stage needs. Everything else is derived.
\end{enumerate}

Units throughout are kilometres and millions of years (Myr); conveniently,
$1\km/\mathrm{Myr} = 1$\,mm/yr.

\section{State}

\subsection{Grids}

Use one lattice of $N$ cells on a sphere of radius $R$, with cell areas $a_i$
and centres $\rr_i$. The lattice must support a constant-time lookup from a
position to the index of the containing cell. A cube-sphere (six square grids
projected onto the sphere) is the simplest choice; HEALPix (Hierarchical Equal
Area isoLatitude Pixelization) gives equal areas.

The same lattice is used in two roles: once fixed in the mantle frame (the
\emph{world grid}), and once per plate, fixed in that plate's co-moving frame
(the \emph{plate grid}). Plate grids are sparse: store only occupied cells, so
total memory is close to $N$ records, not $N$ times the number of plates.

\subsection{Variables}

\begin{table}[h]
\centering
\begin{tabularx}{\textwidth}{@{}llX@{}}
\toprule
Owner & Symbol & Meaning \\
\midrule
Plate $p$ & $q_p$ & Unit quaternion taking plate-frame positions to world positions \\
          & $\om_p$ & Angular velocity (world frame), recomputed every step \\
Plate cell & type & Oceanic or continental \\
          & $h$ & Crust thickness \\
          & $t_0$ & Creation time (age is $t - t_0$) \\
          & $d$ & Damage, $0 \le d \le 1$ \\
          & $t_\mathrm{or}$, kind & Time and kind of last mountain-building or rifting event \\
          & $\hat{\vect{\phi}}$ & Fold direction (unit tangent vector, plate frame) \\
          & $V$ & Accumulated volcanic addition \\
Mantle & $T_{lm}$ & Real spherical harmonic coefficients of temperature anomaly, $1 \le l \le L$ \\
       & hotspots & List of world-frame positions with strength and expiry time \\
Plate pair & polarity & Which plate subducts; time of last convergence \\
\bottomrule
\end{tabularx}
\end{table}

The first four cell fields feed back into the dynamics. The remainder are
history, written for the downstream stage.

\section{One time step}

\begin{enumerate}[leftmargin=*]
\item Update the mantle temperature coefficients (Section~\ref{sec:mantle}).
\item Assemble forces and solve for every plate's angular velocity
(Section~\ref{sec:forces}).
\item Rotate every plate (Section~\ref{sec:kin}).
\item Composite the plates onto the world grid; create crust in gaps and
resolve overlaps (Section~\ref{sec:composite}).
\item Apply within-plate relaxations: erosion, damage healing, hotspot
stamping (Sections~\ref{sec:iso} and~\ref{sec:mantle}).
\item Every $\Delta t_r$: test for welding and rifting
(Sections~\ref{sec:weld} and~\ref{sec:rift}).
\end{enumerate}

\section{Kinematics}\label{sec:kin}

Every rigid motion of a cap on a sphere is a rotation about an axis through
the centre. A point of plate $p$ at world position $\rr$ moves with velocity
\begin{equation}
\vect{v}(\rr) = \om_p \times \rr .
\end{equation}
With $\om_p$ held constant over the step, the orientation update is exact:
\begin{equation}
q_p \leftarrow
\Bigl[\cos\tfrac{\theta}{2},\; \hat{\om}_p \sin\tfrac{\theta}{2}\Bigr] \otimes q_p,
\qquad \theta = |\om_p|\,\Delta t .
\end{equation}
A plate-frame position $\vect{s}$ appears in the world at
$\rr = q_p\,\vect{s}\,q_p^{*}$, and the inverse lookup is
$\vect{s} = q_p^{*}\,\rr\,q_p$.

\section{Compositing: creation and destruction of crust}\label{sec:composite}

For each world cell $i$ and each plate $p$ whose bounding cap contains $\rr_i$,
compute $\vect{s} = q_p^{*}\rr_i q_p$ and look up the plate cell containing
$\vect{s}$. If that cell is occupied, plate $p$ \emph{covers} world cell $i$.

The lookup runs from world to plate, never the reverse. Every world point maps
to exactly one cell of each plate grid, so plate interiors cannot develop
pinholes from lattice aliasing, and crust data is never interpolated while it
moves.

\subsection{Gaps: sea-floor spreading}

A world cell covered by no plate is a gap. Assign gap cells to plates by a
multi-source breadth-first flood from the gap edges, which shares the new
crust roughly evenly between the diverging plates. In the assigned plate,
occupy the cell at $\vect{s}$ with oceanic crust, thickness $h_\mathrm{oc}$,
creation time $t_0 = t$, and zero damage. The line of youngest crust is the
mid-ocean ridge.

\subsection{Overlaps: subduction, accretion, collision}

Define a cell as \emph{subductable} if it is oceanic and thinner than
$h_\mathrm{crit}$. Thickened oceanic crust (island arcs, hotspot plateaus) and
all continental crust are buoyant and are not subductable.

For a world cell covered by plates $p$ and $q$:
\begin{itemize}[leftmargin=*]
\item \textbf{Choose the lower plate.} If the pair already has a recorded
polarity, use it. Otherwise the plate whose cell is subductable goes down; if
both are, the older (denser) one does. Record the polarity on the pair so it
does not flicker from cell to cell.
\item \textbf{Lower cell subductable (subduction).} Delete the cell from the
lower plate. Add its area to the subduction flux $\Phi_i$ at this world cell,
and flag the lower plate's adjacent boundary cells as carrying a slab of the
deleted cell's age.
\item \textbf{Lower cell not subductable (accretion or collision).} Delete
the cell from the lower plate and add its thickness to the upper plate's cell,
conserving crust volume. Set the upper cell's event kind to
\emph{collisional}, $t_\mathrm{or} = t$, raise its damage, and set the fold
direction perpendicular to the relative velocity:
$\hat{\vect{\phi}} \propto \rr \times (\vect{v}_p - \vect{v}_q)$. If the
upper cell is oceanic and the merged thickness exceeds $h_\mathrm{crit}$,
relabel it continental. Record the cell as a contact cell for the pair
(used in Sections~\ref{sec:forces} and~\ref{sec:weld}).
\end{itemize}

Where thickness exceeds $h_\mathrm{max}$, move the excess to neighbouring
cells of the same plate. This spreads a collision zone into a plateau.

\subsection{Volcanic arcs}

On the upper plate, compute the distance $\delta$ from the subduction front by
a breadth-first search limited to a few cells. Thicken the crust at a rate
\begin{equation}
\frac{\partial h}{\partial t} = A\, v_\mathrm{conv}\, g(\delta),
\end{equation}
where $v_\mathrm{conv}$ is the local convergence speed and $g$ is a bump
peaked at $\delta_\mathrm{arc} \approx 200\km$. Add the same amount to $V$ and
raise the damage. When the cell size exceeds $\delta_\mathrm{arc}$ the arc
collapses onto the boundary cell, which is the intended coarse limit. On
oceanic upper plates this builds island arcs, which later become
non-subductable and accrete onto continents.

\section{Isostasy and elevation}\label{sec:iso}

Elevation $E$ is computed on demand from the state. For continental crust,
assume Airy isostasy (the crust floats on the mantle like an iceberg):
\begin{equation}
E = (h - h_0)\left(1 - \frac{\rho_c}{\rho_m}\right),
\end{equation}
with $h_0 \approx 35\km$ the thickness that sits at sea level and
$1 - \rho_c/\rho_m \approx 0.15$. For oceanic crust, depth follows the cooling
of the plate with age $\tau = t - t_0$:
\begin{equation}
E = -\Bigl[z_r + k_z \sqrt{\min(\tau, \tau_\mathrm{flat})}\Bigr],
\end{equation}
with $z_r \approx 2.5\km$, $k_z \approx 0.35\km\,\mathrm{Myr}^{-1/2}$ and
$\tau_\mathrm{flat} \approx 80\Myr$. Oceanic cells thickened beyond
$h_\mathrm{oc}$ receive an additional Airy term.

Erosion is the one downstream process that feeds back, because it thins the
crust. Approximate it as relaxation of excess thickness, with an exact update:
\begin{equation}
h \leftarrow h_0 + (h - h_0)\,e^{-\Delta t/\tau_e} \qquad (h > h_0).
\end{equation}

\section{Mantle temperature field}\label{sec:mantle}

The mantle is represented by a scalar temperature anomaly on the unit sphere,
\begin{equation}
T(\hat{\vect{x}}, t) = \sum_{l=1}^{L}\sum_{m=-l}^{l} T_{lm}(t)\,Y_{lm}(\hat{\vect{x}}),
\end{equation}
using real orthonormal spherical harmonics $Y_{lm}$ and a cutoff $L$ of 4 to
6. The $l = 0$ term is omitted so the mean cannot drift.

\subsection{Evolution}

Each coefficient relaxes linearly and is driven by a source:
\begin{equation}
\frac{dT_{lm}}{dt} = -\lambda_l T_{lm} + S_{lm},
\qquad
\lambda_l = \frac{1 + \epsilon\, l(l+1)}{\tau_m} .
\end{equation}
The $\epsilon$ term is diffusion, which damps short wavelengths faster.
Holding the source constant over a step gives the exact update
\begin{equation}
T_{lm} \leftarrow T_{lm}\,e^{-\lambda_l \Delta t}
 + \frac{S_{lm}}{\lambda_l}\left(1 - e^{-\lambda_l \Delta t}\right).
\end{equation}
The source on the world grid is
\begin{equation}
S_i = -\alpha\,\frac{\Phi_i}{a_i\,\Delta t} + \beta\,C_i ,
\end{equation}
where $\Phi_i$ is the area subducted at cell $i$ during the step and $C_i$ is
1 if the cell is covered by continental crust and 0 otherwise. Subduction
cools the mantle; continents insulate it. Project onto the harmonics with
\begin{equation}
S_{lm} = \sum_i S_i\, Y_{lm}(\hat{\vect{x}}_i)\,\frac{a_i}{R^2} .
\end{equation}

\subsection{Flow}

Surface mantle flow runs from hot to cold:
\begin{equation}
\uu = -k\,\gs T ,
\end{equation}
where $\gs$ is the surface gradient. Evaluate $T$ on the world grid and
difference between neighbours, or differentiate the harmonics analytically.
This flow is purely divergent (poloidal). Rotational (toroidal) surface motion
arises from the rigid plates themselves.

\subsection{Hotspots}

Hotspots are points fixed in the world frame. They are born as a Poisson
process with rate density proportional to $\max(T, 0)$ and live for
$\tau_\mathrm{hs}$. Each step, a hotspot thickens whatever plate cell lies
above it at a fixed rate, adds to $V$, and raises damage. Plate motion turns
this into an age-progressive volcanic chain.

\section{Forces and the torque balance}\label{sec:forces}

Plate inertia is negligible, so at every instant the torques on each plate sum
to zero. All torques are taken about the centre of the sphere.

\subsection{Basal drag}

The traction on the base of plate $p$ at cell $i$ is
$c_i\,(\uu_i - \om_p \times \rr_i)$. The coupling $c_i$ is small beneath
oceanic crust and grows with thickness and age beneath continents, which have
deep roots. Using
$\rr \times (\om \times \rr) = (R^2 \mathsf{I} - \rr\rr^{\mathsf T})\,\om$,
the drag torque is
\begin{equation}
\vect{\tau}^\mathrm{drag}_p = \vect{b}_p - \mathsf{M}_p\,\om_p,
\qquad
\vect{b}_p = \sum_{i \in p} c_i a_i\, \rr_i \times \uu_i,
\qquad
\mathsf{M}_p = \sum_{i \in p} c_i a_i \left(R^2 \mathsf{I} - \rr_i \rr_i^{\mathsf T}\right).
\end{equation}
$\mathsf{M}_p$ has the form of an inertia tensor weighted by drag and is
positive definite for any plate of more than one cell.

\subsection{Slab pull}

On each boundary cell $j$ of plate $p$ flagged as carrying a slab, apply a
force per unit length
\begin{equation}
\vect{f}_j = F_\mathrm{sp}\,\min\!\Bigl(1, \sqrt{\tau_j/\tau_\mathrm{ref}}\Bigr)\,\nn_j ,
\end{equation}
where $\tau_j$ is the age of the crust most recently subducted there and
$\nn_j$ is the unit tangent vector pointing out of the plate across the
boundary. Older slabs are colder and pull harder. The flag expires
$t_\mathrm{slab}$ after convergence stops, which represents the slab breaking
off. With boundary segment length $\ell_j$:
\begin{equation}
\vect{\tau}^\mathrm{sp}_p = \sum_j \ell_j\, \rr_j \times \vect{f}_j .
\end{equation}

\subsection{Topographic body force (optional)}

High ground pushes outward. As a crude proxy for the gradient of
gravitational potential energy, apply a force per unit area
$\vect{f}^g_i = -\gamma\,\gs E$ to every cell. Over young ocean floor this is
ridge push; over thick continental crust it favours extension. Its torque is
$\vect{\tau}^{g}_p = \sum_i a_i\,\rr_i \times \vect{f}^g_i$.

\subsection{Contact resistance}

At each contact cell $j$ between plates $p$ and $q$, apply a force on $p$ of
$-\kappa_j(\om_p - \om_q) \times \rr_j$ and the opposite force on $q$. The
coefficient $\kappa_j$ is large for continental collision and small for
subduction. The resulting torque on $p$ is
$-\mathsf{C}_{pq}(\om_p - \om_q)$ with
\begin{equation}
\mathsf{C}_{pq} = \sum_{j \in \mathrm{contact}(p,q)} \kappa_j
\left(R^2 \mathsf{I} - \rr_j \rr_j^{\mathsf T}\right).
\end{equation}

\subsection{Solve}

Setting the total torque on every plate to zero gives a linear system in the
angular velocities:
\begin{equation}
\Bigl(\mathsf{M}_p + \sum_{q} \mathsf{C}_{pq}\Bigr)\om_p
 - \sum_{q} \mathsf{C}_{pq}\,\om_q
 = \vect{b}_p + \vect{\tau}^\mathrm{sp}_p + \vect{\tau}^{g}_p .
\label{eq:solve}
\end{equation}
For $P$ plates this is a symmetric positive definite system of size $3P$,
solved directly each step. Without contacts it decouples into one
$3\times 3$ solve per plate. Because the contact term is treated implicitly,
colliding plates slow smoothly instead of oscillating.

\section{Welding}\label{sec:weld}

A pair of plates welds when their contact is continental and the mean
relative speed over the contact cells has stayed below $v_\mathrm{weld}$ for
longer than $t_\mathrm{weld}$. To weld $q$ into $p$: copy each occupied cell
of $q$ into $p$'s grid by nearest-cell lookup through
$q_p^{*}\,q_q$, set the damage of the contact cells to 1, mark their event
kind as \emph{suture}, and delete $q$. This is the only occasion on which
crust is resampled.

\section{Rifting}\label{sec:rift}

A plate rifts along a line that is both weak and under tension.

\subsection{Strength}

Each plate cell has a strength (force per unit length)
\begin{equation}
S = S_\mathrm{base}\,(1 - d)\,\exp\!\left[-\frac{\max(T, 0)}{T_w}\right],
\end{equation}
where $T$ is the mantle anomaly beneath the cell. For oceanic crust
$S_\mathrm{base}$ is large and grows with $\sqrt{\tau}$, so oceans almost
never rift. For continental crust $S_\mathrm{base}$ is smaller and grows
slowly with time since the last event, so old undisturbed interiors (cratons)
are strong. Damage is raised by sutures, arcs, hotspots and collisions, and
heals exactly:
\begin{equation}
d \leftarrow d\,e^{-\Delta t/\tau_h}.
\end{equation}

\subsection{Candidate lines}

Every $\Delta t_r$, for each plate, place $K$ seed points evenly around its
perimeter. From each seed run Dijkstra's shortest-path algorithm across the
plate's cells with edge cost $S\,\ell$. Keep the paths between seed pairs
that divide the plate into parts $A$ and $B$ each holding at least a fraction
$f_\mathrm{min}$ of its area.

\subsection{Tension across a candidate line}

The plate is rigid, so it has no strain, but the load a line $\Gamma$ must
carry follows from a free-body argument. Let $\vect{\tau}_A$ be the net
external torque on part $A$ alone, evaluated at the plate's actual $\om_p$:
\begin{equation}
\vect{\tau}_A = \sum_{i \in A} a_i\,\rr_i \times
\Bigl[c_i(\uu_i - \om_p \times \rr_i) + \vect{f}^g_i\Bigr]
+ \bigl(\text{slab pull and contact torques on the boundary of } A\bigr).
\end{equation}
Since the whole plate is in balance, $\vect{\tau}_B = -\vect{\tau}_A$, and
$\Gamma$ must transmit $-\vect{\tau}_A$ to $A$. A uniform line tension
$\sigma$ along $\Gamma$ exerts a torque $-\sigma\,\vect{g}$ on $A$, where
\begin{equation}
\vect{g} = \int_\Gamma \rr \times \nn \, d\ell
\end{equation}
and $\nn$ is the unit tangent vector normal to $\Gamma$ pointing from $B$ into
$A$. Projecting the balance onto this mode gives
\begin{equation}
\sigma = \frac{\vect{\tau}_A \cdot \vect{g}}{\vect{g}\cdot\vect{g}} .
\end{equation}
Positive $\sigma$ means the two parts are being pulled apart. The projection
discards shear and bending along the line.

\subsection{Damage growth and splitting}

Let $\bar S$ be the mean strength along $\Gamma$. Where $\sigma > \bar S$,
grow the damage of the cells on the path:
\begin{equation}
\frac{\partial d}{\partial t} = k_r\,\frac{\sigma - \bar S}{\bar S} .
\end{equation}
Damage lowers both $\bar S$ and the path cost, so a stressed line is selected
again at the next check and weakens faster: the rift localises. When $\bar S$
falls below $\varepsilon\,S_\mathrm{base}$, split the plate. Flood-fill from
either side of $\Gamma$ to label parts $A$ and $B$, move the cells of $B$
into a new plate with the same orientation $q_p$, and mark cells bordering
$\Gamma$ as \emph{rifted margin}. The next solve of Eq.~\eqref{eq:solve}
gives the two parts separate angular velocities, and the gap between them
fills with new ocean floor.

If the tension subsides before the split, the partial damage remains and
heals slowly. This is a failed rift, and a preferred site for the next one.

\section{Initial conditions}

Scatter plate seeds on the sphere and grow them by randomised flood fill to
give irregular plates. Mark a fraction (about 30\%) of the area as continental
crust of thickness $h_0$ in a few contiguous patches, and the rest as oceanic
with random ages. Give the $T_{lm}$ small random values. Run for several
hundred Myr before using the output, so the initial geometry is forgotten.

\section{Behaviour under coarsening}

\begin{table}[h]
\centering
\begin{tabularx}{\textwidth}{@{}lX@{}}
\toprule
Component & Why it survives a larger $\Delta t$ or cell size \\
\midrule
Plate rotation & Exact for any $\Delta t$ \\
Crust transport & No resampling, hence no numerical diffusion \\
Creation and destruction & Set by swept area \\
Mantle field, erosion, healing & Exact exponential updates \\
Torque balance & Implicit linear solve; sums converge as the grid refines \\
Rifting, hotspots & Rate-based \\
Arcs, plateaus & Collapse onto single cells when unresolved \\
\bottomrule
\end{tabularx}
\end{table}

The main limit is that $\om_p$ is frozen during a step. A plate should not
cross more than a few cells per step, or boundary interactions are skipped:
with speeds up to $100\km/\mathrm{Myr}$ and $100\km$ cells, that suggests
$\Delta t$ of 1 to 3 Myr. Transform faults, trenches, back-arc basins and
anything narrower than a cell are left to the next stage.

\section{Output to the next stage}

For each world cell, export: crust type, thickness, age, isostatic elevation,
event kind and time, fold direction, volcanic addition, damage, plate
velocity, and distance to the nearest boundary of each kind. A higher
resolution stage can then add ridges aligned with the fold direction,
roughness that decays with time since the event, trenches along subduction
fronts, individual volcanoes where $V$ is high, and rivers and erosion on top.

\section{Parameters}

Values are starting points of the right order of magnitude. Force-like
parameters ($F_\mathrm{sp}$, $c$, $\kappa$, $\gamma$, $k$, $S_\mathrm{base}$)
only matter through their ratios; tune them so plate speeds fall between 10
and $100\km/\mathrm{Myr}$.

\begin{table}[h]
\centering
\begin{tabular}{@{}lll@{}}
\toprule
Symbol & Meaning & Start value \\
\midrule
$R$ & Planet radius & $6371\km$ \\
$\Delta t$ & Time step & 1--3 Myr \\
$\Delta t_r$ & Interval between weld and rift checks & 10 Myr \\
$h_\mathrm{oc}$, $h_0$, $h_\mathrm{max}$ & Oceanic, reference, maximum thickness & 7, 35, $70\km$ \\
$h_\mathrm{crit}$ & Thickest subductable oceanic crust & $15\km$ \\
$\tau_e$ & Erosion time & 200 Myr \\
$L$ & Harmonic cutoff & 4--6 \\
$\tau_m$ & Mantle relaxation time & 150 Myr \\
$\epsilon$ & Mantle diffusion factor & 0.05 \\
$\tau_\mathrm{hs}$ & Hotspot lifetime & 100 Myr \\
$\tau_\mathrm{ref}$ & Slab age for full pull & 80 Myr \\
$t_\mathrm{slab}$ & Slab persistence after convergence stops & 10 Myr \\
$\delta_\mathrm{arc}$ & Arc distance from trench & $200\km$ \\
$v_\mathrm{weld}$, $t_\mathrm{weld}$ & Welding speed and duration & $5\km/\mathrm{Myr}$, 20 Myr \\
$\tau_h$ & Damage healing time & 1000 Myr \\
$K$, $f_\mathrm{min}$ & Rift seeds per plate, minimum part fraction & 12, 0.1 \\
\bottomrule
\end{tabular}
\end{table}

\section{Status of the claims}\label{sec:status}

\paragraph{Established.} Rigid plate motion as rotation about an Euler pole.
Airy isostasy and the square-root-of-age law for ocean depth~\cite{parsons}.
Slab pull as the dominant driving force, with fast plates being those
attached to slabs and slow plates those carrying continents~\cite{forsyth}.
The observation that rifts tend to reuse old sutures and avoid
cratons~\cite{wilson}.

\paragraph{Hypothesis in the literature.} That supercontinents warm the mantle
beneath them and thereby promote their own breakup~\cite{anderson,gurnis}.
This is contested; here it is adopted because it gives a mechanism cheaply.

\paragraph{Borrowed design.} Rigid plates with per-plate crust and overlap
rules follow the approach of Cortial et al.~\cite{cortial} and of the
Tectonics.js project by Carl Davidson.

\paragraph{Untested construction.} The harmonic mantle field and its source
terms, the gradient flow law, the topographic body force, the coupled contact
solve, and the whole rifting procedure (strength field, path search, projected
tension, damage growth) are assembled for this document and have not been
run. That supercontinent cycles emerge is a conjecture. Expect to tune
thresholds, and expect at least one of these pieces to need redesign.

\paragraph{References.} The citations below are given from memory and should
be checked before being relied upon.

\begin{thebibliography}{9}
\bibitem{cortial} Y.~Cortial, A.~Peytavie, E.~Galin, E.~Gu\'erin,
``Procedural Tectonic Planets,'' \emph{Computer Graphics Forum} 38(2), 2019.
\bibitem{forsyth} D.~Forsyth, S.~Uyeda, ``On the relative importance of the
driving forces of plate motion,'' \emph{Geophysical Journal of the Royal
Astronomical Society} 43, 163--200, 1975.
\bibitem{parsons} B.~Parsons, J.~G.~Sclater, ``An analysis of the variation of
ocean floor bathymetry and heat flow with age,'' \emph{Journal of Geophysical
Research} 82, 803--827, 1977.
\bibitem{wilson} J.~T.~Wilson, ``Did the Atlantic close and then re-open?''
\emph{Nature} 211, 676--681, 1966.
\bibitem{anderson} D.~L.~Anderson, ``Hotspots, polar wander, Mesozoic
convection and the geoid,'' \emph{Nature} 297, 391--393, 1982.
\bibitem{gurnis} M.~Gurnis, ``Large-scale mantle convection and the
aggregation and dispersal of supercontinents,'' \emph{Nature} 332, 695--699,
1988.
\end{thebibliography}

\end{document}
